\section{Related Work}
\label{sec:related_work}

\subsection{Memory-Augmented Language Models}

The challenge of extending LLM context has motivated significant research. Early approaches focused on architectural modifications to transformers for longer sequences \citep{bulatov2022recurrent, liu2023think}. However, these methods still face fundamental limits and computational costs that scale quadratically with context length.

\textbf{MemoryBank} \citep{zhong2024memorybank} introduced a human-inspired memory system with forgetting mechanisms, where memories strengthen when recalled and decay over time. While effective for user modeling, MemoryBank lacks versioning and multi-agent support.

\textbf{MemGPT} \citep{packer2023memgpt} proposed an OS-inspired hierarchical memory with explicit paging between context tiers. MemGPT's function-call-based memory management influenced our design, but MemSt extends this with Git-like versioning and structured knowledge graphs.

\textbf{A-MEM} \citep{xu2025mem} applied Zettelkasten principles of interconnected notes to agent memory, demonstrating that dynamic linking improves multi-hop reasoning. MemSt's Knowledge Graph Extraction v2 builds on this insight with explicit relational modeling and 80+ domain ontologies.

\subsection{Retrieval-Augmented Generation}

RAG systems \citep{lewis2020retrieval} augment LLMs with external knowledge retrieval. While effective for static knowledge bases, standard RAG struggles with conversational memory due to:
\begin{itemize}
    \item Chunking destroying conversational coherence
    \item No distinction between user preferences and general knowledge
    \item Static retrieval without importance weighting
\end{itemize}

MemSt addresses these limitations through tiered memory with lifecycle management and token-budget-aware retrieval.

\subsection{Knowledge Graphs for LLMs}

Several systems incorporate structured knowledge:

\textbf{Zep} \citep{rasmussen2025zep} uses graph representations for entity relationships but suffers from high token overhead ($>$600K tokens per conversation) and slow construction times (hours for background processing).

\textbf{GraphRAG} \citep{edge2024local} extracts entity graphs from document corpora for question answering. MemSt adapts similar extraction techniques but applies them incrementally to streaming conversations with domain-specific ontologies.

\textbf{KGLM} \citep{logan2019baracks} demonstrated that structured knowledge improves language modeling. MemSt extends this to the agent memory domain with runtime extraction and maintenance.

\subsection{Vector Databases}

Systems like Chroma, Milvus, and Pinecone provide efficient vector storage and similarity search. While MemSt uses HNSW (similar to these systems), it adds:
\begin{itemize}
    \item Hybrid search combining vector and keyword signals
    \item Tiered storage with different performance characteristics
    \item Versioning and provenance tracking
\end{itemize}

\subsection{Cloud Memory Services}

\textbf{Mem0} \citep{chhikara2025mem0} provides a managed memory service with impressive accuracy but requires external API dependency. MemSt offers comparable performance with local-first deployment.

\textbf{Letta} (formerly MemGPT) provides hosted agent memory with sleep-time compute but lacks versioning and Git-like operations.

\subsection{Multi-Agent Systems}

Research on multi-agent memory has focused on:
\begin{itemize}
    \item \textbf{Shared memory spaces}: Risk of cross-contamination \citep{wang2024multi}
    \item \textbf{Message passing}: High communication overhead
    \item \textbf{Federated learning}: Privacy-preserving but complex
\end{itemize}

MemSt's worktree approach draws inspiration from Git's branching model, providing isolation without the overhead of full replication.

\subsection{Content-Addressable Storage}

Git's content-addressable model has been applied to other domains:
\begin{itemize}
    \item \textbf{DVC} (Data Version Control) for machine learning datasets
    \item \textbf{IPFS} for distributed web content
    \item \textbf{Hashicorp Vault} for secrets management
\end{itemize}

To our knowledge, MemSt is the first to apply this model comprehensively to LLM agent memory with full versioning semantics.
